% !TeX program = lualatex
% !TeX spellcheck = en_GB
\documentclass[10pt,a4paper]{article}
\usepackage{fontspec}

\usepackage[UKenglish]{babel}
\usepackage{amsmath, amsfonts, amsthm, amssymb, mathtools, stmaryrd}

\usepackage[hyphens]{url}

\usepackage{tikz}
\usepackage{tikz-cd}
\usepackage{tkz-graph}
\usetikzlibrary{arrows,backgrounds,positioning,fit,cd,babel}

\usepackage[pdfusetitle,hidelinks]{hyperref}
\usepackage[capitalise,noabbrev]{cleveref}
\newcommand{\creflastconjunction}{, and\nobreakspace} % https://tex.stackexchange.com/questions/161338/can-cleveref-be-made-to-use-the-oxford-comma-for-multiple-citations


\usepackage{csquotes}
%\usepackage{lmodern}


\usepackage{relsize}
\usepackage{microtype}
% \usepackage{pdfpages}
%\usepackage[a4paper,margin=1in]{geometry}

\usepackage{thm-restate}
\usepackage{subcaption}
\usepackage{enumitem}
\usepackage{comment}



\usepackage{todonotes}


\linespread{1.06}


\usepackage[style=numeric-comp,maxbibnames=99,backend=biber,bibencoding=utf8,sorting=nty]{biblatex} % ,sorting=none,backref=true


\usepackage{booktabs}


\addbibresource{literature.bib}
\AtEveryBibitem{
	\clearfield{urldate}
	\clearfield{urlyear}	
	\clearfield{urlmonth}
	\clearfield{issn}	
	\clearfield{month}
	\clearfield{day}		
	\clearfield{language}
	\clearfield{Language}
	\clearfield{langid}
	\clearfield{note}
}
%https://tex.stackexchange.com/questions/424774/print-url-only-if-doi-not-present
\DeclareSourcemap{
	\maps[datatype=bibtex]{
		\map[overwrite]{
			\step[fieldsource=doi, final]
			\step[fieldset=url, null]
			\step[fieldset=eprint, null]
			\step[fieldset=isbn, null]
		}  
	}
}


\theoremstyle{definition}
\newtheorem{definition}{Definition}[section]
\newtheorem{notation}[definition]{Notation}
\newtheorem{example}[definition]{Example}
\newtheorem{remark}[definition]{Remark}

\theoremstyle{plain}
\newtheorem{lemma}[definition]{Lemma}
\newtheorem{corollary}[definition]{Corollary}
\newtheorem{theorem}[definition]{Theorem}
\newtheorem{conjecture}[definition]{Conjecture}
\newtheorem{question}[definition]{Question}
\newtheorem{proposition}[definition]{Proposition}
\newtheorem{observation}[definition]{Observation}
\newtheorem{fact}[definition]{Fact}
\Crefname{fact}{Fact}{Facts}
\newtheorem{proviso}[definition]{Proviso}

\theoremstyle{remark}
\newtheorem{claim}{Claim}[definition]
\renewcommand*{\theclaim}{\thetheorem\alph{claim}}
\Crefname{claim}{Claim}{Claims}
\newenvironment{claimproof}[1][Proof of Claim]{\begin{proof}[#1] \renewcommand{\qedsymbol}{$\lrcorner$}}{ \end{proof}}
\setlist[enumerate, 1]{font=\upshape, noitemsep, nolistsep}
\setlist[enumerate, 2]{font=\upshape, noitemsep, nolistsep}
\setlist[itemize, 1]{noitemsep, nolistsep,font=\upshape}
\setlist[itemize, 2]{noitemsep, nolistsep,font=\upshape}

\DeclarePairedDelimiter{\norm}{\lVert}{\rVert}
\DeclarePairedDelimiter{\abs}{\lvert}{\rvert}

\newcommand{\cart}{\mathbin\square}

\newcommand{\Aa}{{\cal A}}
\newcommand{\Bb}{{\cal B}}
\newcommand{\Cc}{{\cal C}}
\newcommand{\Dd}{{\cal D}}
\newcommand{\Ee}{{\cal E}}
\newcommand{\Ff}{{\cal F}}
\newcommand{\Gg}{{\cal G}}
\newcommand{\Hh}{{\cal H}}
\newcommand{\Ii}{{\cal I}}
\newcommand{\Jj}{{\cal J}}
\newcommand{\Kk}{{\cal K}}
\newcommand{\Ll}{{\cal L}}
\newcommand{\Mm}{{\cal M}}
\newcommand{\Nn}{{\cal N}}
\newcommand{\Oo}{{\cal O}}
\newcommand{\Pp}{{\cal P}}
\newcommand{\Qq}{{\cal Q}}
\newcommand{\Rr}{{\cal R}}
\newcommand{\Ss}{{\cal S}}
\newcommand{\Tt}{{\cal T}}
\newcommand{\Uu}{{\cal U}}
\newcommand{\Vv}{{\cal V}}
\newcommand{\Ww}{{\cal W}}
\newcommand{\Xx}{{\cal X}}
\newcommand{\Zz}{{\cal Z}}


\DeclareMathOperator{\tw}{tw}
\DeclareMathOperator{\hdtw}{hdtw}
\DeclareMathOperator{\mn}{mn}
\DeclareMathOperator{\sub}{sub}
\DeclareMathOperator{\emb}{emb}
\DeclareMathOperator{\perm}{perm}
\DeclareMathOperator{\ml}{ml}
\DeclareMathOperator{\hdim}{hdim}
\DeclareMathOperator{\vc}{vc}
\DeclareMathOperator{\mln}{mln}
\DeclareMathOperator{\tr}{tr}
\DeclareMathOperator{\aut}{aut}
\DeclareMathOperator{\Sub}{Sub}
\DeclareMathOperator{\cl}{cl}
\DeclareMathOperator{\id}{id}
\DeclareMathOperator{\cc}{lvc}



\tikzset{
	vertex/.style={draw,circle,fill=gray},
	every node/.style={anchor=center},
	lbl/.style={color=lightgray}
}

\title{Symmetric Algebraic Circuits and \\ Homomorphism Polynomials\footnote{This work was presented at the \emph{17th Innovations in Theoretical Computer Science Conference} (\textsmaller{ITCS} 2026) \cite{conference_version}.}}
\author{Anuj Dawar \orcidlink{0000-0003-4014-8248} \and Benedikt Pago \orcidlink{0000-0001-6377-1230} \and Tim Seppelt \orcidlink{0000-0002-6447-0568}}
% \author{Anonymous authors}

%\date{}

\renewcommand{\phi}{\varphi}
\renewcommand{\epsilon}{\varepsilon}

\newcommand{\Aut}{\mathbf{Aut}}
\newcommand{\Sym}{\mathbf{Sym}}
\newcommand{\Alt}{\mathbf{Alt}}
\newcommand{\Stab}{\mathbf{Stab}}
\newcommand{\StabP}{\Stab^{\bullet}}
\newcommand{\Orb}{\mathbf{Orb}}

\DeclareMathOperator{\child}{children}




\newcommand{\bbN}{\mathbb{N}}
\newcommand{\bbQ}{\mathbb{Q}}
\newcommand{\bbF}{\mathbb{F}}




\definecolor{lightgray}{rgb}{0.60, 0.60, 0.61} % lipicsLightGray
\definecolor{gray}{rgb}{0.31, 0.31, 0.33} % lipicsBulletGray
\definecolor{yellow}{rgb}{0.99, 0.78, 0.07}



\let\sup\relax
\DeclareMathOperator{\sup}{sup}

\DeclareMathOperator{\imm}{imm}
\DeclareMathOperator{\sgn}{sgn}
\DeclareMathOperator{\Match}{\# Match}
\DeclareMathOperator{\maxOrb}{maxOrb}
\DeclareMathOperator{\maxSup}{maxSup}



\newcommand{\VP}{\mathsf{VP}}
\newcommand{\VNP}{\mathsf{VNP}}
\newcommand{\VFPT}{\mathsf{VFPT}}
\newcommand{\VW}{\mathsf{VW}[1]}

\newcommand{\FPT}{\mathsf{FPT}}
\newcommand{\sharpW}{\#\mathsf{W}[1]}


\usepackage[margin=22mm]{geometry}
